\chapter{The Model, and the Object of Study}\label{ch:model}
\label{ch:background}\label{ch:graphs}\label{ch:graphical}\label{ch:development}

This chapter fixes what the thesis computes and the language it computes in.
A Markov chain generates a clean trajectory; an Ornstein--Uhlenbeck channel
corrupts every frame independently; the object of study is the joint score of
the corrupted sequence. Two justifications come before any analysis. The first
is that the quantity worth studying is the joint score of a whole trajectory
and not a per-frame score; the second is that the stationary AR(1) chain is a convenient
model in which to study it.

\section{Markov sequences}\label{sec:stationarity}\label{sec:markov}

We briefly recall the notions of stationarity, Markov dependence, and
Toeplitz covariance used throughout, together with the one point among them
that is easy to get wrong and matters later \citep{brockwell1991time}.

A process $(X_k)_{k \in \mathbb{Z}}$ is \emph{strictly stationary} if its
finite-dimensional laws are invariant under shifts in time: for every
$n \in \mathbb{N}$, every $k_1 < \dots < k_n$, and every $h \in \mathbb{Z}$,
\begin{equation}
  \big(X_{k_1+h}, \dots, X_{k_n+h}\big)
  \;\overset{d}{=}\;
  \big(X_{k_1}, \dots, X_{k_n}\big).
  \label{eq:strict-stationary}
\end{equation}
A process with finite second moments is \emph{weakly}, or
\emph{covariance-}, \emph{stationary} if its mean is constant,
\begin{equation}
  \E[X_k] = \mu ,
\end{equation}
and its covariance depends only on the distance between two indices,
\begin{equation}
  \operatorname{Cov}(X_j,X_k) = \gamma(j-k).
  \label{eq:weak-stationary}
\end{equation}
Strict stationarity implies weak stationarity whenever second moments exist.
For a Gaussian process the converse also holds, because a Gaussian law is
determined completely by its mean and covariance.

Weak stationarity alone gives the covariance matrix of any segment
$(X_0,\ldots,X_{L-1})$ of length $L$ a \emph{Toeplitz} form, constant along
every diagonal:
\begin{equation}
  \Sigma_{ij} = \gamma(i-j), \qquad i,j = 0,\ldots,L-1,
  \qquad
  \Sigma =
  \begin{pmatrix}
    \gamma(0) & \gamma(1) & \gamma(2) & \cdots \\
    \gamma(1) & \gamma(0) & \gamma(1) & \cdots \\
    \gamma(2) & \gamma(1) & \gamma(0) & \cdots \\
    \vdots    & \vdots    & \vdots    & \ddots
  \end{pmatrix},
  \label{eq:toeplitz}
\end{equation}
using $\gamma(-d) = \gamma(d)$ for a real-valued process.

\begin{definition}[Markov chain]
A sequence $(X_k)_{k \geq 0}$ on a state space $\mathcal{X}$ is a \emph{Markov
chain} if for every $k$ and every measurable $A \subseteq \mathcal{X}$,
\begin{equation}
  \Prob\big(X_{k+1} \in A \,\big|\, X_k, X_{k-1}, \dots, X_0\big)
  \;=\; \Prob\big(X_{k+1} \in A \,\big|\, X_k\big) .
  \label{eq:markov-property}
\end{equation}
It is \emph{time-homogeneous} if a single \emph{transition kernel}
$\kernel(x' \,|\, x)$, with $\int \kernel(x'|x)\,\dd x' = 1$ for every $x$,
governs every step; every chain in this thesis is of this kind.
\end{definition}

The Markov property \eqref{eq:markov-property} is exactly a statement about
\emph{factorisation}: the joint law of the first $L$ states is
\begin{equation}
  p(x_0, x_1, \dots, x_{L-1})
  \;=\; p_0(x_0)\, \prod_{k=0}^{L-2} \kernel(x_{k+1} \,|\, x_k) .
  \label{eq:chain-factorisation}
\end{equation}
This is the most heavily used equation of the thesis. It says the joint
distribution of a Markov trajectory is a product of \emph{local} terms, each
touching two consecutive variables: the structure that makes the chain a tree
in the factor-graph sense of Chapter~\ref{ch:graphs}, that makes its precision
matrix tridiagonal in the Gaussian case of Chapter~\ref{ch:gaussian}, and that
makes belief propagation exact. Marginals evolve by integrating the kernel,
$p_{k+1}(x') = \int \kernel(x'|x)\, p_k(x)\,\dd x$, and a law $\pi$ with
$\pi = \kernel \pi$ is \emph{stationary}: a chain started in $\pi$ remains in
it, and its trajectory is then strictly stationary.

\begin{example}[The Gaussian AR(1) chain]\label{ex:ar1}
The running example throughout is the scalar Gaussian autoregression of order
one,
\begin{equation}
  a_{k+1} \;=\; \alpha\, a_k + \eta_k ,
  \qquad
  \eta_k \stackrel{\text{iid}}{\sim} \Nd(0, \ivar),
  \qquad k = 0, \dots, L-2 ,
  \label{eq:ar1}
\end{equation}
with kernel $\kernel(a' \mid a) = \Nd(a'; \alpha a, \ivar)$. Taking
$\ivar = 1 - \alpha^2$ and $a_0 \sim \Nd(0,1)$ gives $a_k \sim \Nd(0,1)$ for
every $k$ and
\begin{equation}
  \operatorname{Cov}(a_j,a_k) = \alpha^{|j-k|}, \qquad
  \Sigma =
  \begin{pmatrix}
    1 & \alpha & \alpha^2 & \cdots & \alpha^{L-1} \\
    \alpha & 1 & \alpha & \cdots & \alpha^{L-2} \\
    \alpha^2 & \alpha & 1 & \cdots & \alpha^{L-3} \\
    \vdots & \vdots & \vdots & \ddots & \vdots \\
    \alpha^{L-1} & \alpha^{L-2} & \alpha^{L-3} & \cdots & 1
  \end{pmatrix},
  \label{eq:ar1-toeplitz}
\end{equation}
Toeplitz by \eqref{eq:toeplitz}. Because $a_0$ is chosen to be the fixed point
of $\kernel$ itself --- $\Nd(0,1)$ solves $\pi = \kernel \pi$ exactly when the
innovation is $\Nd(0,\ivar)$ --- the chain is started in its own stationary
law, hence strictly stationary by the general fact above, and not merely
weakly stationary by the covariance computation. In this Gaussian baseline the
three notions coincide completely: the same model is a time-homogeneous
Markov chain, a strictly stationary process, and a weakly stationary process
with Toeplitz covariance, all at once. Every chapter from
Chapter~\ref{ch:gaussian} onwards concerns \eqref{eq:ar1}, one of its
non-Gaussian variants, or what an estimator can recover of it through noise.
\end{example}

The coincidence is special to the Gaussian case, and breaks the moment
Chapter~\ref{ch:laplace} keeps $a_0 \sim \Nd(0,1)$ and $\alpha$ fixed but
replaces $\eta_k$ with a Laplace innovation of matched variance $1-\alpha^2$.
Only the weak notion is available for the
non-Gaussian chains this thesis studies.


\section{The Ornstein--Uhlenbeck corruption channel}\label{sec:ou}

The noising process used by diffusion models, and by every result of this
thesis, is the \emph{Ornstein--Uhlenbeck process}
\citep{uhlenbeck1930theory}: Brownian motion pulled back to the origin by a
linear spring,
\begin{equation}
  \dd X_t \;=\; -X_t\,\dd t \;+\; \sqrt{2}\;\dd W_t ,
  \qquad X_0 = a .
  \label{eq:ou-sde}
\end{equation}
The drift and diffusion constants are the variance-preserving normalisation of
the diffusion literature \citep{song2021scorebased}; any other choice rescales
time. The process is the canonical stationary Gaussian Markov diffusion: by Doob's
characterisation \citep{doob1942brownian} it is, up to affine rescaling of
state and time, the only one-dimensional stationary Gaussian Markov process
with continuous sample paths and nondegenerate stationary variance. Stationary
there describes the process started from its invariant law $\Nd(0,1)$; the run
from the fixed $X_0 = a$ of \eqref{eq:ou-sde} merely relaxes to it. It is exactly solvable: the
derivation is Toolbox~\ref{tb:ou-solution}.

\begin{toolbox}[tb:ou-solution]{Solving the Ornstein--Uhlenbeck SDE by an
integrating factor}
\emph{Integrating factor.} Multiply \eqref{eq:ou-sde} by $e^{t}$ and apply
the product rule to $e^{t}X_t$ --- ordinary, not Itô's, because $e^t$ carries
no stochastic part and contributes no quadratic-variation term:
\begin{equation*}
  \dd\big(e^{t}X_t\big)
  \;=\; e^{t}X_t\,\dd t + e^{t}\,\dd X_t
  \;=\; e^{t}X_t\,\dd t + e^{t}\big(-X_t\,\dd t + \sqrt2\,\dd W_t\big)
  \;=\; \sqrt2\, e^{t}\,\dd W_t .
\end{equation*}
The drift term cancels exactly, leaving a pure stochastic differential.

\emph{Integration.} Integrating from $0$ to $t$ and using $X_0 = a$,
\begin{equation*}
  e^{t}X_t - a \;=\; \sqrt2 \int_0^t e^{s}\,\dd W_s
  \qquad\Longrightarrow\qquad
  X_t \;=\; a\,e^{-t} \;+\; \sqrt2\, e^{-t}\int_0^t e^{s}\,\dd W_s .
\end{equation*}

\emph{The stochastic integral is Gaussian.} Write
$I_t := \sqrt2\,e^{-t}\int_0^t e^{s}\,\dd W_s$. This is an Itô integral of a
\emph{deterministic} integrand against Brownian motion: it is the $L^2$ limit
of Riemann-sum increments $\sum_i \sqrt2\,e^{-t}e^{s_i}
\big(W_{s_{i+1}}-W_{s_i}\big)$, each term Gaussian and independent of the
others, so the limit $I_t$ is Gaussian too, with mean zero by linearity of
the expectation. Its variance follows from the Itô isometry,
\begin{equation*}
\begin{aligned}
  \operatorname{Var}(I_t)
  &= \E\bigg[\Big(\sqrt2\,e^{-t}\int_0^t e^{s}\,\dd W_s\Big)^{\!2}\bigg]
   = 2e^{-2t}\int_0^t e^{2s}\,\dd s \\
  &= 2e^{-2t}\cdot\frac{e^{2t}-1}{2}
   = 1-e^{-2t}
   = \Deltat .
\end{aligned}
\end{equation*}

\emph{Conclusion.} $X_t = a\,e^{-t} + I_t$ with $I_t \sim \Nd(0,\Deltat)$
independent of $a$, so
\begin{equation*}
  X_t \mid X_0 = a \;\sim\; \Nd\big(a\,e^{-t},\, \Deltat\big),
  \qquad \Deltat = 1 - e^{-2t},
\end{equation*}
which is \eqref{eq:ou-channel} below. \hfill$\square$
\end{toolbox}

Corrupting a data point $a$ for a time $t$ therefore
produces
\begin{equation}
  x \;=\; a\, e^{-t} \;+\; \sqrt{\Deltat}\;\varepsilon,
  \qquad \varepsilon \sim \Nd(0, 1):
  \label{eq:ou-channel}
\end{equation}
the signal is \emph{contracted} by $e^{-t}$ while isotropic noise of variance
$\Deltat$ fills in. At $t = 0$ the channel is the identity; as $t \to \infty$,
$x \sim \Nd(0,1)$ regardless of $a$. The signal-to-noise ratio
$e^{-2t}/\Deltat$ sweeps continuously from $\infty$ to $0$ and diverges as
$1/2t$ for small $t$. The channel interpolates between the data distribution and a
standard Gaussian, which is what a diffusion model needs.


\section{How the object of study was arrived at}\label{sec:dev-history}

The estimand of this thesis was not chosen at the outset. It was reached by
getting it wrong first; the correction is a consequential step in the project.

\paragraph{The propagator hypothesis.} The work began from a structural guess
about sequences. If a trajectory is generated by a fixed dynamical rule, the score at one frame
ought to be obtainable from the score at the previous one by a fixed operator.
Writing $S_{:,u}$ for the block of the score at internal time $u$, the
conjecture was that some \emph{propagator} $P_t$ satisfies
\begin{equation}
  S_{:,u}(x) \;\approx\; P_t\big(S_{:,u-1}(x)\big).
  \label{eq:dev-propagator}
\end{equation}
If such an operator existed and were cheap, generation over a long sequence
would cost one solve plus $L-1$ applications of $P_t$. The programme adopted to
test \eqref{eq:dev-propagator} was deliberately austere: exactly solvable
models only, closed forms wherever possible, so that the focus is on understanding

\paragraph{Six toy models, and a mistake they shared.} A sequence of solvable
objects was built to probe \eqref{eq:dev-propagator}: a static Gaussian
mixture under the channel; a two-frame joint law; a two-dimensional rotating
process $a_{u+1} = R_{\Delta\theta}\, a_u + \xi_u$; a billiard in a box; a
scalar AR(1) chain examined on the whole $(u,t)$ plane; and a diffusion on the
circle. The rotating model and the circle are the direct ancestors of
Chapter~\ref{ch:ring}, and the AR(1) study is the ancestor of
Chapters~\ref{ch:gaussian} and~\ref{ch:gaussian-bp}.

All but one of them computed the wrong object. What was evaluated was the
\emph{per-frame marginal} score
\begin{equation}
  s(x_u, u, t) \;=\; \nabla_{x_u} \log p_t(x_u \mid u),
  \label{eq:dev-marginal-score}
\end{equation}
which integrates out every other frame. Doing so is not a mild approximation.
Driving generation with the stacked marginal scores is operationally equivalent
to replacing the joint score by the score of the product law
$\prod_u p_0(a_u \mid u)$ --- the model one would have if the clean law
factorised over frames, which it does not for any correlated process and, in
particular, for no model in the list.
The quantity being studied had the frames' dependence removed from it before
the study began, so no propagator could have been found in it: there was
nothing left for $P_t$ to propagate.

\paragraph{The correction.} The object was corrected to the \emph{joint} score of the whole trajectory,
$S(x,t) = \nabla_{(x_0,\dots,x_{L-1})} \log P_t(x_0,\dots,x_{L-1})$, whose
$k$-th block is the posterior expectation
\begin{equation}
  S_k(x,t) \;=\; \frac{e^{-t}\,\E\big[a_k \,\big|\, x_0,\dots,x_{L-1}\big] - x_k}{\Deltat},
  \label{eq:dev-joint-block}
\end{equation}
taken over the full posterior of the clean trajectory given all $L$
observations.

\section{Why the joint score}\label{sec:dev-correction}

A dynamic datum is a sequence
$x=(x_0,\ldots,x_{L-1})$, so its reverse diffusion is driven by the
\emph{joint} score
\begin{equation}
  s_t(x)
  =
  \nabla_x \log p_t(x_0,\ldots,x_{L-1}),
  \label{eq:dev-joint-score}
\end{equation}
rather than by the collection of one-frame marginal scores
\[
  s^{\mathrm{marg}}_{t,k}(x_k)
  =
  \partial_{x_k}\log p_{t,k}(x_k).
\]
The distinction is structural:
\begin{equation}
  \boxed{
  s_{t,k}(x_0,\ldots,x_{L-1})
  \text{ may depend on the other frames, whereas }
  s^{\mathrm{marg}}_{t,k}(x_k)
  \text{ cannot.}}
  \label{eq:dev-joint-vs-marginal}
\end{equation}

The smallest example already contains the whole argument. Consider two
consecutive states of the stationary Gaussian AR(1) chain,
\begin{equation}
  \begin{pmatrix}a_0\\a_1\end{pmatrix}
  \sim
  \Nd\!\left(
    0,
    \Sigma_0
  \right),
  \qquad
  \Sigma_0
  =
  \begin{pmatrix}
    1 & \alpha\\
    \alpha & 1
  \end{pmatrix},
\end{equation}
and corrupt both coordinates independently through
\begin{equation}
  x_i=e^{-t}a_i+\sqrt{\Deltat}\,z_i,
  \qquad
  \Deltat=1-e^{-2t},
  \qquad
  z_i\stackrel{\mathrm{iid}}{\sim}\Nd(0,1).
\end{equation}
Since the added noise is independent,
\begin{align}
  \operatorname{Var}(x_i)
  &=e^{-2t}\operatorname{Var}(a_i)+\Deltat
    =e^{-2t}+1-e^{-2t}=1,\\
  \operatorname{Cov}(x_0,x_1)
  &=e^{-2t}\operatorname{Cov}(a_0,a_1)
    =\alpha e^{-2t}.
\end{align}
Writing
\begin{equation}
  r_t:=\alpha e^{-2t},
\end{equation}
the noisy pair is therefore
\begin{equation}
  \begin{pmatrix}x_0\\x_1\end{pmatrix}
  \sim
  \Nd\!\left(
    0,
    \Sigma_t
  \right),
  \qquad
  \Sigma_t=
  \begin{pmatrix}
    1&r_t\\
    r_t&1
  \end{pmatrix}.
  \label{eq:dev-L2-cov}
\end{equation}

The two marginals are thus
\begin{equation}
  x_0,x_1\sim\Nd(0,1),
  \qquad
  s^{\mathrm{marg}}_{t,0}(x_0)=-x_0,
  \qquad
  s^{\mathrm{marg}}_{t,1}(x_1)=-x_1,
  \label{eq:dev-L2-marginal}
\end{equation}
with no $\alpha$ anywhere.

The joint law is different. Since
\begin{equation}
  \Sigma_t^{-1}
  =
  \frac{1}{1-r_t^2}
  \begin{pmatrix}
    1&-r_t\\
    -r_t&1
  \end{pmatrix},
\end{equation}
its density is
\begin{equation}
  p_t(x_0,x_1)
  =
  \frac{1}{2\pi\sqrt{1-r_t^2}}
  \exp\!\left[
    -\frac{x_0^2-2r_tx_0x_1+x_1^2}
           {2(1-r_t^2)}
  \right].
  \label{eq:dev-L2-density}
\end{equation}
Differentiating its logarithm gives
\begin{align}
  s_{t,0}(x_0,x_1)
  &=
  -\frac{x_0-r_tx_1}{1-r_t^2},\\
  s_{t,1}(x_0,x_1)
  &=
  -\frac{x_1-r_tx_0}{1-r_t^2},
\end{align}
or
\begin{equation}
  \boxed{
  s_t(x_0,x_1)
  =
  -\frac{1}{1-r_t^2}
  \begin{pmatrix}
    x_0-r_tx_1\\
    x_1-r_tx_0
  \end{pmatrix}.}
  \label{eq:dev-L2-score}
\end{equation}
This is the missing information: the marginal score of frame $0$ is
$-x_0$, whereas its joint score contains the observation $x_1$ through
$r_tx_1$. The temporal coupling has disappeared from the marginal law but
remains explicitly in the joint score.

For a length-$L$ stationary Gaussian AR(1) chain the same calculation becomes
\begin{equation}
  \boxed{
  \Sigma_t
  =
  e^{-2t}\Sigma_0+\Deltat I,
  \qquad
  (\Sigma_0)_{ij}=\alpha^{|i-j|},
  \qquad
  s_t(x)=-\Sigma_t^{-1}x.}
  \label{eq:dev-general-gaussian}
\end{equation}
Hence
\begin{equation}
  (\Sigma_t)_{ij}
  =
  \begin{cases}
    1, & i=j,\\[2mm]
    e^{-2t}\alpha^{|i-j|}, & i\neq j.
  \end{cases}
  \label{eq:dev-general-cov}
\end{equation}
The diagonal, and therefore every one-frame Gaussian marginal, is independent
of $\alpha$; all information about the temporal coupling sits in the
off-diagonal structure.

This also makes the Fisher-information statement immediate. A one-frame
density is simply
\[
  p_{t,k}(x_k;\alpha)
  =
  \frac{1}{\sqrt{2\pi}}e^{-x_k^2/2},
\]
which does not contain $\alpha$. Therefore
\begin{equation}
  \partial_\alpha\log p_{t,k}(x_k;\alpha)=0
  \qquad\Longrightarrow\qquad
  \boxed{
  I_k(\alpha)
  =
  \E\!\left[
    \bigl(\partial_\alpha\log p_{t,k}(X_k;\alpha)\bigr)^2
  \right]
  =0.}
  \label{eq:dev-fisher-zero}
\end{equation}
Hence, its Fisher information is exactly zero.

\begin{remark}[What changes outside the Gaussian case]
\label{rmk:dev-nongaussian}
The conclusion that the marginal is always $\Nd(0,1)$ is special to Gaussian
data. For a general scalar input $A$,
\begin{equation}
  X_t=e^{-t}A+\sqrt{\Deltat}\,Z
  \quad\Longrightarrow\quad
  \varphi_{X_t}(u)
  =
  \varphi_A(e^{-t}u)
  \exp\!\left(-\frac{\Deltat u^2}{2}\right).
  \label{eq:dev-cf}
\end{equation}
If $A\sim\Nd(0,1)$, the two Gaussian factors combine to give
$\varphi_{X_t}(u)=e^{-u^2/2}$ exactly. By contrast, for a unit-variance
Laplace variable,
$\varphi_A(u)=(1+u^2/2)^{-1}$, and therefore
\begin{equation}
  \varphi_{X_t}(u)
  =
  \frac{
    \exp[-\Deltat u^2/2]
  }{
    1+e^{-2t}u^2/2
  },
\end{equation}
which is not Gaussian for any finite $t$. Gaussian corruption therefore
preserves the Gaussian family exactly, while for non-Gaussian data it only
drives the marginal progressively towards Gaussianity as $t$ increases.
\end{remark}

The lesson is therefore not that marginal scores are always trivial, but that
\emph{one-frame marginals do not determine dependence}. Figure~\ref{fig:toymodel}
shows precisely this contrast.

\begin{figure}[t]
  \centering
  \includegraphics[width=\linewidth]{fig_marginal_blindness.pdf}
  \caption[Same marginals, different joint score]{
  \textbf{Same one-frame marginals, different joint score.}
  In both cases $x_0,x_1\sim\Nd(0,1)$, so a frame-wise model sees the same
  marginal densities and the same scores $(-x_0,-x_1)$.
  When $\alpha=0$ the joint density is circular and the joint score is indeed
  radial. When $\alpha>0$, the off-diagonal covariance
  $r_t=\alpha e^{-2t}$ tilts the joint density and changes each score
  component through the other coordinate, as in
  \eqref{eq:dev-L2-score}. The dependence is therefore visible in the joint
  score and invisible in either marginal.
  }
  \label{fig:toymodel}
\end{figure}

Finally, by the score--posterior identity \eqref{eq:dev-joint-block},
\begin{equation}
  s_{t,k}(x)
  =
  \frac{
    e^{-t}\E[a_k\mid x_0,\ldots,x_{L-1}]-x_k
  }{\Deltat},
\end{equation}
so computing the joint score is an inference problem over the whole
trajectory. This is the point at which the graphical structure of the
Markov chain becomes useful, and motivates the belief-propagation
construction developed below.

\section{Markov random fields and factor graphs}\label{sec:mrf}
\label{sec:factor-graphs}

Graphical models assign structure to distributions the way a precision matrix
assigns it to a Gaussian, and for arbitrary variables. A collection
$x = (x_i)_{i \in V}$ on an undirected graph $G$ is a \emph{Markov random field}
if each variable is conditionally independent of the rest given its neighbours,
$p(x_i \mid x_{V \setminus \{i\}}) = p(x_i \mid x_{\partial i})$. For a
strictly positive density the Hammersley--Clifford theorem
\citep{besag1974spatial} makes that equivalent to factorisation over the cliques
of $G$,
\begin{equation}
  p(x) \;=\; \frac{1}{Z} \prod_{C \in \mathrm{cliques}(G)} \psi_C(x_C)
  \;=\; \frac{1}{Z}\, \exp\Big[ -\sum_C E_C(x_C) \Big] ,
  \label{eq:hammersley-clifford}
\end{equation}
so conditional independence, graphical locality and additive local energies are
one thing seen three ways. For Gaussian fields it becomes matrix-concrete: the
energy is $\tfrac12 a^\T J a - h^\T a$, and the graph is exactly the sparsity
pattern of $J$. A Gaussian Markov chain therefore has a tridiagonal precision,
which Chapter~\ref{ch:gaussian} derives and uses throughout.

If each variable is observed through
its own independent channel, the likelihood factors touch one variable each and
no clique is enlarged: \emph{the posterior of an MRF under independent per-node
observations is an MRF on the same graph}. Independent unary observation
changes the potentials, never the graph; correlated observation noise would add
edges.

An MRF records which variables interact but not how the joint factorises: a
triangle of pairwise terms and one triple term have the same graph. The
\emph{factor graph} makes the factorisation itself graphical
\citep{kschischang2001factor}. Given
$p(x) = \frac{1}{Z}\prod_a \psi_a(x_{\partial a})$, it is the bipartite graph
with a variable node per $x_i$, a factor node per $\psi_a$, and an edge whenever
$x_i \in x_{\partial a}$. A connected factor graph with one fewer edge than
nodes is a \emph{tree}, and trees are the tractable case: on them marginals and
partition functions are exact under the recursion of the next section, whatever
the alphabet. The factor graph of a Markov chain observed through per-node
noise is a tree of a particularly simple shape, a spine of variables and
transition factors with one evidence leaf per variable, as
Chapter~\ref{ch:gaussian} verifies by counting.
Figure~\ref{fig:rbm-graph}
draws the bipartite structure used in RBMs as an example.

\begin{figure}[t]\centering
\begin{tikzpicture}[
  latent/.style={circle, draw, minimum size=7mm, inner sep=1pt},
  obs/.style={circle, draw, fill=black!10, minimum size=7mm, inner sep=1pt},
  node distance=9mm, scale=0.92, every node/.style={transform shape}]
\foreach \a in {1,2,3,4}{
  \node[latent] (h\a) at (1.6*\a, 1.6) {$h_{\a}$};
}
\foreach \i in {1,...,6}{
  \node[obs] (v\i) at (1.1*\i, 0) {$v_{\i}$};
}
\foreach \a in {1,2,3,4}{
  \foreach \i in {1,...,6}{
    \draw[black!45, thin] (h\a) -- (v\i);
  }
}
\node[left=5mm of h1] {Hidden layer};
\node[left=5mm of v1] {Visible layer};
\end{tikzpicture}
\caption[The RBM as a bipartite factor graph]{The restricted Boltzmann
machine's graph: every hidden unit connects to every visible unit, and no
edge runs within a layer. Because the energy has only
visible--hidden cliques, conditioning on one layer leaves the other in
disjoint unary cliques, so $p(h\mid v)=\prod_a p(h_a\mid v)$ and
$p(v\mid h)=\prod_i p(v_i\mid h)$: the same conditional-independence
argument as the general MRF case above, specialised to a graph with no
within-layer edges at all.}
\label{fig:rbm-graph}
\end{figure}

\section{Belief Propagation}\label{sec:sum-product}

On a factor graph, sum--product belief propagation (BP) passes messages
along every edge \citep{pearl1988probabilistic, kschischang2001factor,
mezard2009information}. Variable-to-factor messages multiply the
incoming factor messages from all \emph{other} neighbours, and
factor-to-variable messages integrate the factor against the incoming
variable messages:
\begin{equation}
  \nu_{i \to a}(x_i) \propto
    \prod_{b \in \partial i \setminus a} \hat\nu_{b \to i}(x_i),
  \qquad
  \hat\nu_{a \to i}(x_i) \propto
    \int \psi_a(x_{\partial a})
    \prod_{j \in \partial a \setminus i} \nu_{j \to a}(x_j)\;
    \dd x_{\partial a \setminus i},
  \label{eq:bp-rules}
\end{equation}
with the belief at node $i$ proportional to
$\prod_{a \in \partial i} \hat\nu_{a \to i}(x_i)$. The ``exclude the
recipient'' rule is the heart of the algorithm: the message $i \to a$
summarises everything node $i$ knows \emph{except} what it learned from
$a$, so that no piece of evidence is ever counted twice along a path.

Figure~\ref{fig:bp-messages} draws the two rules of \eqref{eq:bp-rules}.

\begin{figure}[t]\centering
\includegraphics[width=\linewidth]{fig_mezard_bp_messages.png}
\caption[The two sum--product update rules]{The two halves of
\eqref{eq:bp-rules}, drawn on the local neighbourhood each one uses.
\emph{Left:} the message into $j$ that will be sent on to $a$, a function of
the incoming factor messages $\hat\nu_{b\to j}$ from every other factor
neighbour $b \neq a$. \emph{Right:} the message out of $a$ towards $j$, a
function of the incoming variable messages $\nu_{k\to a}$ from every other
variable neighbour $k \neq j$. In the notation of this thesis, $j$ here is
$i$ and $k$ is $j$; the original also carries iteration superscripts
$(t)$/$(t{+}1)$; because this thesis's graphs are trees
(Theorem~\ref{thm:tree-exact} below), one outward sweep and one inward sweep
of each rule is exact, and no iteration index is needed here. Reproduced
from \citet{mezard2009information}, their Figure~14.3.}
\label{fig:bp-messages}
\end{figure}

\begin{theorem}[Tree exactness of sum--product]\label{thm:tree-exact}
On a tree-shaped factor graph, one inward and one outward pass of
\eqref{eq:bp-rules} compute the exact node and factor marginals of $p$
\citep[Theorem~14.1]{mezard2009information}. The resulting beliefs are
unique; the individual unnormalised messages are defined only up to positive
multiplicative constants.
\end{theorem}

\emph{Proof. (Not formal)} Each edge of a tree splits the graph into two disjoint halves, so
the message crossing it summarises exactly the half behind it and no evidence
reaches a node twice. The recursion terminates with every marginal having
absorbed each factor once. \hfill$\square$

 On graphs with cycles the same updates,
iterated, define \emph{loopy} BP, which is extraordinarily useful but
approximate \citep{weiss2001correctness, malioutov2006walk}; none of that
is needed in this thesis, because our graphs are trees.

\begin{figure}[t]\centering
\begin{tikzpicture}[
  latent/.style={circle, draw, minimum size=7mm, inner sep=1pt},
  obs/.style={circle, draw, fill=black!10, minimum size=7mm, inner sep=1pt},
  fac/.style={rectangle, draw, fill=black, minimum size=2.2mm, inner sep=0pt},
  node distance=9mm, scale=0.92, every node/.style={transform shape}]
\foreach \i in {1,2,3,4}{
  \node[latent] (a\i) at (2.1*\i, 0) {$a_{\i}$};
  \node[fac] (f\i) at (2.1*\i, -1.0) {};
  \node[obs] (x\i) at (2.1*\i, -2.0) {$x_{\i}$};
  \draw (a\i) -- (f\i) -- (x\i);
}
\foreach \i/\j in {1/2,2/3,3/4}{
  \node[fac] (w\i) at (2.1*\i+1.05, 0) {};
  \draw (a\i) -- (w\i) -- (a\j);
}
\node[above=1.5mm of w1] {\small $\kernel(a_{k}\mid a_{k-1})$};
\node[right=2mm of f4] {\small $\like(x_k\mid a_k)$};
\draw[-{Stealth}, thick] (1.15,1.35) -- (2.55,1.35) node[midway,above] {\small $\lmsg$};
\draw[-{Stealth}, thick] (9.35,1.35) -- (7.95,1.35) node[midway,above] {\small $\rmsg$};
\end{tikzpicture}
\caption[Posterior factor graph of a noised Markov chain]{The posterior factor
graph of a Markov chain observed through the coordinatewise channel. Open
circles are latent variables, shaded circles observations, squares factors. Each
transition factor $\kernel(a_k \mid a_{k-1})$ joins two adjacent latent
variables; each observation enters through a \emph{unary} factor attached to a
single latent variable, because the channel corrupts every coordinate
independently. Conditioning therefore reweights node potentials without creating
any edge among the $a_k$, so the posterior graph is still a chain and
Theorem~\ref{thm:tree-exact} applies to it. Arrows show the two message sweeps
of \eqref{eq:bp-rules}. The construction fails as soon as the forward noise is
correlated across coordinates: the likelihood then touches several $a_k$ at
once and the graph acquires edges.}
\label{fig:posterior-factor-graph}
\end{figure}

\paragraph{The representation problem.} For \emph{discrete} variables
each message is a finite probability vector and each update in
\eqref{eq:bp-rules} is finite arithmetic. For \emph{continuous} variables
each message is a probability \emph{density}, an infinite-dimensional
object, and the update integrals must be carried out in function space.
Tree exactness (a property of the graph) says nothing about
representability of the messages (a property of the model's function
class); keeping the two apart is essential, and the distinction is where
all the difficulty of the non-Gaussian research chapter lives. There are
three ways to make continuous BP computable: restrict to models whose
messages close in a finite-dimensional family (the Gaussian route, next
section); represent messages numerically on a grid (the route our
non-Gaussian reference solver takes, Chapter~\ref{ch:laplace}); or
project each message onto a tractable family and accept the error (the
route whose cost this thesis measures).
